\subsection{Subsets and Power Sets}

\begin{method} \label{meth:proof-by-containment} \hlt{Proof by Two Containments}\\
To prove $A=B$, prove both inclusions.
\begin{enumerate}[label=(\roman*)]
\setlength{\itemsep}{0pt}
\item To prove $A\subseteq B$, let $x\in A$ be arbitrary and use the definitions and hypotheses to show $x\in B$.
\item To prove $B\subseteq A$, let $x\in B$ be arbitrary and use the definitions and hypotheses to show $x\in A$.
\item Once both containments are proved, conclude $A=B$ by extensionality of sets.
\end{enumerate}
If both sides are described by membership conditions, this is a chain, $x\in A \Longleftrightarrow \cdots \Longleftrightarrow x\in B$.
\end{method}

\hlt{Placeholder}: Add extensionality examples, proper subset notation, power-set examples, and common containment proof patterns.
